\section{Bootstrap 重建分析流程的不确定性}
\label{sec:13-Machine-Learning/08-Sampling-and-Inference/04}

一份院内风险模型的报告放在你桌上：从八百名患者的记录里筛出十二个变量，交叉验证调好正则强度，最终模型的 AUC 是 \(0.83\)。评审只问了一句——换一批患者重做一遍，这个 \(0.83\) 会变成多少。谁也答不上来。前三节处理的都是给定概率模型之后算不出的积分，这里的困境不同：整条分析流程里有筛选、有调参、有插补，它根本写不成一个能求期望的解析式，谈不上后验，也谈不上抽样分布。

Bootstrap 的做法是把``重新收集一份数据''这件事本身模拟出来。记完整分析流程为 \(A\)，原数据为 \(D\)，估计量 \(\widehat\theta=A(D)\)。从经验分布（或拟合好的模型）生成复样 \(D^{*b}\)，重新跑一遍流程，

\[
\widehat\theta^{*b}=A(D^{*b}),\qquad b=1,\ldots,B,
\]

复样结果的经验分布就用来估标准误、偏差、区间和选择稳定性。原理与统计部分 \hyperref[sec:11-Mathematical-Statistics/06-Resampling-and-Model-Selection/01]{``Bootstrap 近似估计量的抽样波动''} 里讲的一致，这里关心的是它在一条真实分析流水线上怎样用错。

先把一件事钉死：\(B\) 是模拟次数，不是新数据。加大 \(B\) 只让 bootstrap 自身的计算误差变小，原始的八百名患者所含的信息一分未增。\(B=10000\) 的区间比 \(B=200\) 的稳定，但两者围绕的宽度是同一个，取决于样本量和流程的脆弱程度。

\subsection{抽样单位由数据怎样产生决定，不由表格形状决定}

如果个体确实独立同分布，按行有放回抽即可。可这份数据里同一名患者有多次住院记录，独立单位是患者而非记录行。逐行抽样把同一患者的相关观测拆散，等于假装样本量比实际大，区间会明显偏窄。时间序列的相邻点相关，要用分块 bootstrap 保住短程依赖；分层抽样设计要在层内复样；回归的固定设计下可以改抽残差，异方差时用 wild bootstrap。

动手之前先画一遍数据的产生层级：谁被随机抽进样本，哪些观测共享同一个患者、科室、批次或时间冲击。复样方案要模仿的是这套抽样机制，不是数据在数据库里的存储格式。这一步搞错，后面所有区间都不必看了。

\subsection{数据自适应的步骤留在循环外，区间就会假窄}

报告里那十二个变量是从八十个候选里筛出来的，正则强度是交叉验证选的。如果每个复样只对这十二个变量、这个固定的正则强度重新拟合，那么筛选和调参带来的波动一点都没进入区间——而它们往往是最大的一部分。换一批患者，被选中的很可能是另外十二个变量。

\begin{center}
\begin{tikzpicture}[>=stealth]
\draw[->] (-4.4,0) -- (4.6,0) node[right]{\(\widehat\theta^{*}\)};
\draw[very thick] plot[smooth] coordinates
  {(-4.0,0.04) (-3.2,0.16) (-2.4,0.42) (-1.6,0.80) (-0.8,1.10)
   (0,1.22) (0.8,1.10) (1.6,0.80) (2.4,0.42) (3.2,0.16) (4.0,0.04)};
\node at (2.9,1.15) {整条流程重跑};
\draw[dashed,thick] plot[smooth] coordinates
  {(-1.5,0.03) (-1.1,0.22) (-0.8,0.62) (-0.5,1.20) (-0.2,1.85)
   (0,2.05) (0.2,1.85) (0.5,1.20) (0.8,0.62) (1.1,0.22) (1.5,0.03)};
\node at (-2.6,1.9) {只重拟合最后一步};
\draw[gray] (-2.7,1.75) -- (-0.9,1.55);
\draw[|-|,very thick] (-2.9,-0.55) -- (2.9,-0.55);
\node[right] at (3.0,-0.55) {诚实的区间};
\draw[|-|,very thick,dashed] (-1.0,-1.05) -- (1.0,-1.05);
\node[right] at (3.0,-1.05) {假窄的区间};
\draw[gray,dashed] (0,-1.3) -- (0,2.25);
\end{tikzpicture}
\end{center}

\emph{把筛选与调参锁在循环外面，复样之间就少了一大半波动，区间随之缩水。}

图里那条虚线分布之所以窄，不是因为估计更精确，而是因为每个复样都被喂了同一份从原数据里偷看来的变量清单。正确的 \(A(D^{*b})\) 要在每个复样内部从头做：重新拟合缺失值插补器与标准化参数，重新做特征筛选，用复样内部的验证重新选超参数，再拟合最终模型，最后算出目标系数、预测或性能指标。这比跑一次贵上百倍，但只有它才对应``重新做一遍这项研究''。同样的道理在 \hyperref[sec:11-Mathematical-Statistics/06-Resampling-and-Model-Selection/03]{``交叉验证评估的是完整训练流程''} 里已经出现过一次：被评估的对象是流程，不是流程的最后一环。

如果目标是部署性能而非某个系数，还要为每个复样留出未参与训练的那部分——袋外样本或独立的外部集——在上面报告指标，否则复样内的乐观偏差会顶着 bootstrap 的名义混进结论。

\subsection{区间公式有适用条件，失败的复样是证据}

最省事的分位数区间直接取 \(\widehat\theta^{*}\) 的经验分位数，它在估计量近似无偏、分布对称时还行，遇上偏差、偏度或参数贴着边界就会失准。基本区间、学生化区间、BCa 各有各的修正方向，换个公式不等于换来普遍正确。Bootstrap 相合性大体要求统计量是经验分布的某种光滑泛函；极值统计量、非正则参数、模型选择之后那种带跳变的统计量、以及只有寥寥几个聚类的情形，都可能踩破这个条件。区间覆盖率属于程序而非某一次估计，这层区分见 \hyperref[sec:11-Mathematical-Statistics/04-Interval-Estimation/01]{``置信度属于程序，而非这次参数''}。

复样里的拟合失败不能悄悄丢掉。某些复样出现完全分离、设计矩阵奇异、某个类别一个样本都没抽到，这不是脏数据，是流程在样本轻微扰动下就不稳定的直接证据。丢掉失败的复样等于只统计顺利的那些，区间会朝好的方向偏。该做的是记录失败率与失败原因，并检查结论对复样方案、随机种子和 \(B\) 的敏感程度。

\subsection{三种不确定性算的是三本不同的账}

Bootstrap 分布近似的是估计量在重复抽样下的波动；Bayesian 后验刻画的是给定这份已观测数据、在所选模型下参数的不确定性；Monte Carlo 标准误衡量的是有限次模拟对某个确定目标的计算误差。三者的数值有时接近，含义不能互换。第一节坚持把 Monte Carlo 标准误写进报告、第二节用有效样本量把相关链折算成独立样本，都是为了不让计算误差混进统计不确定性那本账。

放回那份风险模型报告里：bootstrap 区间回答的是``换一批患者，AUC 会怎样变''；MCMC 后验区间回答的是``就这批患者、就这个模型，参数还剩多少不确定''；把 \(B\) 从 \(200\) 加到 \(10000\)、把链从一万步跑到一百万步，只让各自的数值更稳，患者数还是八百。这三本账都建立在同一个前提上——每个复样都能被完整地重跑一遍。下一节要处理的场景连这个前提也没有：数据一个接一个到达，你必须在每个时刻就给出当前状态的分布，没有回头重算的机会。
